\section{Thảo luận và hạn chế}\label{sec:discussion}

\subsection{Trả lời câu hỏi nghiên cứu}

\pending{Trả lời RQ3 đến RQ6 theo đúng kết cục đọc ở mục~6 (được ủng hộ, trái với giả thuyết, chưa kết
luận được), không mạnh hơn số liệu; nêu rõ tầng vừa và tầng lớn chỉ là bằng chứng sơ bộ.}

\subsection{Đối chiếu với tài liệu}

\pending{Đặt kết quả cạnh Hill (không thấy đặc tả giảm lỗi), Agarwal (đánh đổi tin cậy và chi phí), Mock
và Russo (tác tử thêm độ phức tạp), Jiang (lập kế hoạch ở mức hàm) và kết quả hồi cứu RQ2.}

\subsection{Mối đe dọa đến tính hợp lệ}\label{sec:threats}

\textit{Tính hợp lệ nội.} Đề bài do AI soạn sau khi đã biết lời giải và được An duyệt gộp; ở tầng vừa và
tầng lớn, mục giao diện bắt buộc nêu nhiều tên, nên đề càng gần một bản đặc tả (rõ nhất ở
\texttt{m01}) thì nhánh theo đặc tả càng ít điều để thêm. Test ẩn được viết cho bản cài đặt gốc, vốn được
làm theo đặc tả, nên có thể thiên vị nhánh này. Can thiệp là một gói (thiết kế, kế hoạch, viết test
trước, chạy cả bộ test), nên không tách được phần của tài liệu đặc tả. Tác tử có công cụ web và mạng mở;
\texttt{upstream\_urls} chỉ bắt đường dẫn hiện ra trong nhật ký, nên là cận dưới của việc đọc mã gốc, và
phân tích độ nhạy chỉ bỏ được những lượt bị bắt. Từ khi tầng vừa và tầng lớn bắt đầu, ba lệnh chạy cùng lúc
với tổng tối đa 24 phiên thay vì tám, nên thời gian chạy và số lượt hết giờ trong khoảng đó có thể cao
hơn; mục~\ref{sec:results} so số lượt hết giờ trong và ngoài khoảng chạy chồng. Việc làm mù khi gán nhãn
chỉ là một phần, và nhãn sau phân xử của An không mù.

\textit{Tính hợp lệ cấu trúc.} ``Đạt'' là đạt test ẩn không mang nhãn \texttt{integration}; repository
Postgres, migration và lệnh quản trị chỉ có test \texttt{integration} nên không được chấm. Ở
\NumGradedMediumMinTasks{}, test được chấm chỉ chạy tới \NumGradedMediumMin{} dòng mã của bản vá gốc, nên
theo \texttt{task\_pass} hai nhiệm vụ này gần tầng nhỏ hơn số dòng gợi ý; phân tích độ nhạy với
\texttt{graded\_src\_lines} kiểm điều này. Các tầng khác nhau không chỉ ở số dòng mà còn ở số test, giới
hạn thời gian và độ chi tiết của đề, nên H5 là bằng chứng về ``nhiệm vụ cỡ tính năng'', không về riêng số
dòng. Ngưỡng 35 tên của tầng lớn không có đơn vị đếm cố định và được áp bằng phán đoán; F09 bị loại theo
quyết định có trước chạy thử, không theo ngưỡng này. Chi phí bằng tiền là số quy đổi theo giá API, không
phải số tiền đã trả; với \texttt{codex-luna} chỉ có cận dưới và cận trên.

Lượng token khác nhau rất xa giữa các công cụ. Ở chạy thử, trung vị token vào mỗi lượt của
\texttt{agy-gemini} là \NumPilotAgyInputMedianM{} triệu với \NumPilotAgyTurnsMedian{} lượt gọi mô hình
(\NumPilotAgyRuns{} lượt có nhật ký, gồm cả lượt bị ngắt vì hạn mức), gấp khoảng \NumPilotAgyInputRatio{}
lần mức \NumPilotSonnetInputMedianM{} triệu với \NumPilotSonnetTurnsMedian{} lượt gọi của
\texttt{claude-sonnet}; \NumPilotAgyCacheReadPct{}\% token vào của \texttt{agy-gemini} là token đọc cache.
Token đọc cache rẻ khi quy đổi theo giá API nhưng vẫn bị tính vào hạn mức thuê bao; điều này phù hợp với
việc cả bốn tài khoản Antigravity hết hạn mức tuần trước khi nhóm này chạy xong. Chạy thử dùng cấu hình cũ trên hai nhiệm
vụ, nên đây chỉ là số mô tả, không phải kết quả.

\textit{Tính hợp lệ ngoại.} Một dự án, một ngôn ngữ (Python), \NumTasksMain{} nhiệm vụ backend, và tác
giả nghiên cứu cũng là tác giả dự án. Ba công cụ dùng mô hình khác hãng và khác hạng. Nhánh theo đặc tả
là bản tự động, không có bước người duyệt đặc tả như khi người làm nghề dùng quy trình này. Một số nhiệm
vụ ở các tầng khác nhau chung commit hoặc chung file test ẩn.

\textit{Tính hợp lệ kết luận.} Tầng vừa có \NumMediumTasks{} và tầng lớn có \NumFeatureTasks{} nhiệm vụ,
nên khoảng tin cậy ở hai tầng này chỉ có giá trị tham khảo. Các phép thử phụ không được hiệu chỉnh so sánh
nhiều lần. Tại thời điểm viết, độ phủ thật của khoảng bootstrap cụm cho mô hình hỗn hợp chưa được đo bằng mô
phỏng. Nếu \texttt{agy-gemini} chưa đủ ở hạn chót, RQ4 chỉ còn hai công cụ cho phép thử có khoảng tin cậy.
\pending{số liệu cụ thể của các mục trên sau khi mở kết quả}
